
\chapter{2018年上海卷（秋）}

\section{填空题}

\begin{problem}
行列式 \(\begin{vmatrix}4 & 1\\2 & 5\end{vmatrix}\) 的值为\fillinblank{}.
\end{problem}

\begin{problem}
双曲线 \(\frac{x^2}{4}-y^2=1\) 的渐近线方程为\fillinblank{}.
\end{problem}

\begin{problem}
在 \((1+x)^7\) 的二项展开式中，\(x^2\) 项的系数为\fillinblank{}. （结果用数值表示）
\end{problem}

\begin{problem}
设常数 \(a\in\R\)，函数 \(f(x)=\log_2(x+a)\). 若 \(f(x)\) 的反函数的图像经过点 \((3,1)\)，则 \(a=\)\fillinblank{}.
\end{problem}

\begin{problem}
已知复数 \(z\) 满足 \((1+\i)z=1-7\i\)（\(\i\) 是虚数单位），则 \(\lvert z\rvert=\)\fillinblank{}.
\end{problem}

\begin{problem}
记等差数列 \(\{a_n\}\) 的前 \(n\) 项和为 \(S_n\)，若 \(a_3=0\)，\(a_6+a_7=14\)，则 \(S_7=\)\fillinblank{}.
\end{problem}

\begin{problem}
已知 \(\alpha\in\{-2,-1,-\frac12,\frac12,1,2,3\}\). 若幂函数 \(f(x)=x^\alpha\) 为奇函数，且在 \((0,+\infty)\) 上递减，则 \(\alpha=\)\fillinblank{}.
\end{problem}

\begin{problem}
在平面直角坐标系中，已知点 \(A(-1,0)\)，\(B(2,0)\)，\(E\)，\(F\) 是 \(y\) 轴上的两个动点，且 \(\lvert\overrightarrow{EF}\rvert=2\)，则 \(\overrightarrow{AE}\cdot\overrightarrow{BF}\) 的最小值为\fillinblank{}.
\end{problem}

\begin{problem}
有编号互不相同的五个砝码，其中 \(5\) 克，\(3\) 克，\(1\) 克砝码各一个，\(2\) 克砝码两个. 从中随机选取三个，则这三个砝码的总质量为 \(9\) 克的概率是\fillinblank{}. （结果用最简分数表示）
\end{problem}

\begin{problem}
设等比数列 \(\{a_n\}\) 的通项公式为 \(a_n=q^{n-1}\)（\(n\in\N^*\)），前 \(n\) 项和为 \(S_n\). 若 \(\displaystyle\lim_{n\to\infty}\frac{S_n}{a_{n+1}}=\frac12\)，则 \(q=\)\fillinblank{}.
\end{problem}

\begin{problem}
已知常数 \(a>0\)，函数 \(f(x)=\frac{2^x}{2^x+ax}\) 的图像经过点 \(P\left(p,\frac65\right)\)，\(Q\left(q,-\frac15\right)\). 若 \(2^{p+q}=36pq\)，则 \(a=\)\fillinblank{}.
\end{problem}

\begin{problem}
已知实数 \(x_1\)，\(x_2\)，\(y_1\)，\(y_2\) 满足：\(x_1^2+y_1^2=1\)，\(x_2^2+y_2^2=1\)，\(x_1x_2+y_1y_2=\frac12\)，则 \(\frac{\lvert x_1+y_1-1\rvert}{\sqrt2}+\frac{\lvert x_2+y_2-1\rvert}{\sqrt2}\) 的最大值为\fillinblank{}.
\end{problem}

\section{单选题}

\begin{problem}
设 \(P\) 是椭圆 \(\frac{x^2}{5}+\frac{y^2}{3}=1\) 上的动点，则 \(P\) 到该椭圆的两个焦点的距离之和为（\quad）

\choices
    {\(2\sqrt2\)}
    {\(2\sqrt3\)}
    {\(2\sqrt5\)}
    {\(4\sqrt2\)}
\end{problem}

\begin{problem}
已知 \(a\in\R\)，则“\(a>1\)”是“\(\frac1a<1\)”的（\quad）

\choices
    {充分非必要条件}
    {必要非充分条件}
    {充要条件}
    {既非充分又非必要条件}
\end{problem}

\begin{problem}
《九章算术》中，称底面为矩形而有一侧棱垂直于底面的四棱锥为阳马. 设 \(AA_1\) 是正六棱柱的一条侧棱，如图，若阳马以该正六棱柱的顶点为顶点，以 \(AA_1\) 为底面矩形的一边，则这样的阳马的个数是（\quad）

\bitmapfigure[max width=6.0cm]{img/2018/shanghai/q15_fig1.png}

\choices
    {\(4\)}
    {\(8\)}
    {\(12\)}
    {\(16\)}
\end{problem}

\begin{problem}
设 \(D\) 是含数 \(1\) 的有限实数集，\(f(x)\) 是定义在 \(D\) 上的函数. 若 \(f(x)\) 的图像绕原点逆时针旋转 \(\frac\pi6\) 后与原图像重合，则在以下各项中，\(f(1)\) 的可能取值只能是（\quad）

\choices
    {\(\sqrt3\)}
    {\(\frac{\sqrt3}{2}\)}
    {\(\frac{\sqrt3}{3}\)}
    {\(0\)}
\end{problem}

\section{解答题}

\begin{problem}
已知圆锥的顶点为 \(P\)，底面圆心为 \(O\)，半径为 \(2\).

\begin{enumerate}
\item 设圆锥的母线长为 \(4\)，求圆锥的体积；
\item 设 \(PO=4\)，\(OA\)，\(OB\) 是底面半径，且 \(\angle AOB=90^\circ\)，\(M\) 为线段 \(AB\) 的中点，如图，求异面直线 \(PM\) 与 \(OB\) 所成的角的大小.
\end{enumerate}

\bitmapfigure[max width=6.0cm]{img/2018/shanghai/q17_fig1.png}
\end{problem}

\begin{problem}
设常数 \(a\in\R\)，函数 \(f(x)=a\sin2x+2\cos^2x\).

\begin{enumerate}
\item 若 \(f(x)\) 为偶函数，求 \(a\) 的值；
\item 若 \(f\left(\frac\pi4\right)=\sqrt3+1\)，求方程 \(f(x)=1-\sqrt2\) 在区间 \([-\pi,\pi]\) 上的解.
\end{enumerate}
\end{problem}

\begin{problem}
某群体的人均通勤时间，是指单日内该群体中成员从居住地到工作地的平均用时. 某地上班族 \(S\) 中的成员仅以自驾或公交方式通勤. 分析显示：当 \(S\) 中 \(x\%\)（\(0<x<100\)）的成员自驾时，自驾群体的人均通勤时间 \(f(x)\)（单位：分钟）为

\examdisplaycases{f(x)=}{
30, & 0<x\le30,\\
2x+\frac{1800}{x}-90, & 30<x<100
}

而公交群体的人均通勤时间不受 \(x\) 影响，恒为 \(40\) 分钟. 试根据上述分析结果回答下列问题：

\begin{enumerate}
\item 当 \(x\) 在什么范围内时，公交群体的人均通勤时间少于自驾群体的人均通勤时间？
\item 求该地上班族 \(S\) 的人均通勤时间 \(g(x)\) 的表达式；讨论 \(g(x)\) 的单调性，并说明其实际意义.
\end{enumerate}
\end{problem}

\begin{problem}
设常数 \(t>2\). 在平面直角坐标系 \(xOy\) 中，已知点 \(F(2,0)\)，直线 \(l:x=t\)，曲线 \(\Gamma:y^2=8x\)（\(0\le x\le t\)，\(y\ge0\)），\(l\) 与 \(x\) 轴交于点 \(A\)，与 \(\Gamma\) 交于点 \(B\). \(P\)，\(Q\) 分别是曲线 \(\Gamma\) 与线段 \(AB\) 上的动点.

\begin{enumerate}
\item 用 \(t\) 表示点 \(B\) 到点 \(F\) 的距离；
\item 设 \(t=3\)，\(\lvert FQ\rvert=2\)，线段 \(OQ\) 的中点在直线 \(FP\) 上，求 \(\triangle AQP\) 的面积；
\item 设 \(t=8\)，是否存在以 \(FP\)，\(FQ\) 为邻边的矩形 \(FPEQ\)，使得点 \(E\) 在 \(\Gamma\) 上？若存在，求点 \(P\) 的坐标；若不存在，说明理由.
\end{enumerate}
\end{problem}

\begin{problem}
给定无穷数列 \(\{a_n\}\)，若无穷数列 \(\{b_n\}\) 满足：对任意 \(n\in\N^*\)，都有 \(\lvert b_n-a_n\rvert\le1\)，则称 \(\{b_n\}\) 与 \(\{a_n\}\)“接近”.

\begin{enumerate}
\item 设 \(\{a_n\}\) 是首项为 \(1\)，公比为 \(\frac12\) 的等比数列，\(b_n=a_{n+1}+1\)，\(n\in\N^*\). 判断数列 \(\{b_n\}\) 是否与 \(\{a_n\}\) 接近，并说明理由；
\item 设数列 \(\{a_n\}\) 的前四项为：\(a_1=1\)，\(a_2=2\)，\(a_3=4\)，\(a_4=8\)，\(\{b_n\}\) 是一个与 \(\{a_n\}\) 接近的数列，记集合 \(M=\{x\mid x=b_i,\ i=1,2,3,4\}\)，求 \(M\) 中元素的个数 \(m\)；
\item 已知 \(\{a_n\}\) 是公差为 \(d\) 的等差数列. 若存在数列 \(\{b_n\}\) 满足：\(\{b_n\}\) 与 \(\{a_n\}\) 接近，且在 \(b_2-b_1\)，\(b_3-b_2\)，\ldots，\(b_{201}-b_{200}\) 中至少有 \(100\) 个为正数，求 \(d\) 的取值范围.
\end{enumerate}
\end{problem}
